\section{一次不定方程：特解与正整数解范围}\label{knowledge-diophantine-bounds}\index{Bezout identity}\index{二元一次不定方程}
实现见第~\ref{compact-linear_equation}~节（第~\pageref{compact-linear_equation}~页）。
对不全为零的 $a,b$，令 $g=\gcd(|a|,|b|)>0$。
递归扩展欧几里得先求 $|a|u+|b|v=g$；当且仅当 $g\mid c$ 时有整数解。
把系数乘以 $c/g$，再恢复 $a,b$ 的符号，就得到 $ax_0+by_0=c$。
两个特解相减，再利用 $a/g,b/g$ 互素，可知所有解恰为
\[
x=x_0+\frac bg t,\qquad y=y_0-\frac ag t,\qquad t\in\mathbb Z.
\]
若 $a=b=0$，则 $c=0$ 时所有整数对都是解，否则无解；此时不能使用除以 $g$ 的通解式。

底层 \texttt{exgcd} 的本库调用约定为非负系数，上限 $2^{63}$；
返回非负的 $g$，$(0,0)$ 约定返回 0。
包装函数先转为 128 位再取相反数，因而可以处理 \texttt{LLONG\_MIN}。
缩放后的特解也用 128 位保存；不要提前转回 \texttt{long long}。
即使特解可表示，直接验证 $ax+by$ 的两个乘积仍可能超过 128 位，本地验证使用任意精度整数。

对 P5656 的正系数 $a,b,c$，设 $d_x=b/g,d_y=a/g$，
先将 $x_0$ 调整到 $x_{\min}\in[1,d_x]$，并计算与它配对的
$y_{\max}=(c-ax_{\min})/b$。
若 $y_{\max}>0$，正整数解个数为
\[
N=\left\lfloor\frac{y_{\max}-1}{d_y}\right\rfloor+1,
\quad x_{\max}=x_{\min}+(N-1)d_x,
\quad y_{\min}=y_{\max}-(N-1)d_y.
\]
否则不存在同时为正的解，但 $x$ 和 $y$ 仍各自有最小正值；
分别按周期归一化即可，它们通常不是同一组解。正值归一化使用
\texttt{(v \% d + d - 1) \% d + 1}，不要把余数 0 当作正整数。
驱动中 $a,b,c\le10^9$，归一化后的乘积不超过 $10^{18}$，故该部分可用 64 位。

双版本以小范围有符号输入穷举、任意精度等式证书、64 位端点和斐波那契递归链验证。
完整 P5656 驱动另验证 200000 组输入，包含正解数量与四个极值、无解和无正解分支；
普通与 ASan/UBSan 测试均通过，在线 AC 待补。

来源：kuangbin 2.3；WIDA 数论打印稿的 exgcd 与不定方程例题。

\url{https://oi-wiki.org/math/number-theory/bezouts/}。

底层接口见第~\ref{compact-extended_gcd}~节（第~\pageref{compact-extended_gcd}~页）。

\newpage
\subsection{有符号矩形范围：把两维约束化成一个参数区间}
现在允许 $a,b,c$ 为负或零，要求
\[
L_x\le x\le R_x,\qquad L_y\le y\le R_y.
\]
先处理空区间；任意一维下界大于上界，答案为0。
若 $a=b=0$，$c\ne0$ 时无解；$c=0$ 时答案为
$(R_x-L_x+1)(R_y-L_y+1)$，不能套用一维参数族。
否则先检查 $g\mid c$，再取得上面的特解与步长
$d_x=b/g$、$d_y=-a/g$。不要把负系数或负步长直接取绝对值而保留旧特解。

对任一坐标 $z=z_0+dt$，条件 $L\le z\le R$ 给出：
\[
\begin{array}{c|c}
 d>0 & \lceil(L-z_0)/d\rceil\le t\le\lfloor(R-z_0)/d\rfloor\\[2pt]
 d<0,\ k=-d>0 & \lceil(z_0-R)/k\rceil\le t\le\lfloor(z_0-L)/k\rfloor\\[2pt]
 d=0 & L\le z_0\le R\text{ 时不限制 }t\text{，否则无解}
\end{array}
\]
例如 $6x-9y=3$，$x\in[-5,4],y\in[-1,5]$ 时，解恰为 $(2,1)$ 和 $(-1,-1)$。
分别计算 $x,y$ 允许的区间，取交集 $[T_l,T_r]$；解数是
$\max(0,T_r-T_l+1)$。不全为零的系数保证至少一个步长非零，
有限矩形因此会给出有限参数区间；实现时用“尚无上下界”标记，不必猜一个足够大的哨兵。
通解对整数 $t$ 是单射，且已包含全部整数解，所以交集计数没有遗漏或重复。
若要实际列举答案，再依次代入这些整数 $t$，输出成本为 $O(\text{解数})$。

这里的 floor/ceil 是数学取整，不能用 C++ 的负数向零截断替代。
对正分母 $D$，令 C++ 商余为 $q=n/D,r=n\%D$，则
\[
\lfloor n/D\rfloor=q-[r<0],\qquad
\lceil n/D\rceil=q+[r>0].
\]
例如 $-5/2$ 的 floor 是 $-3$、ceil 是 $-2$。
不要使用只对非负数成立的 \texttt{(n+D-1)/D}，也不要在未证明安全时先计算 \texttt{-n}。
参数平移 $x'_0=x_0+d_xs,y'_0=y_0+d_ys$ 只会把允许的 $t$ 整体平移 $-s$，答案不变；
这是检查实现符号的一个独立不变量。

范围运算比求一个特解要求更强的整数宽度：即使输入端点都是64位，差值和区间长度也未必能用64位保存。
当两系数全零、矩形覆盖全部有符号64位坐标时，答案为 $2^{128}$，连无符号128位也放不下。
本节通式按精确整数理解；实现前证明每一步范围，不能证明时使用任意精度整数。
尤其不要用窄类型计算长度后才转宽，也不能在等式证书中先让 $ax,by$ 溢出再相消。
求特解之后只需常数次整数运算；大整数版本的位复杂度取决于操作数长度，不能据此承诺固定字长常数耗时。

本节的有界小整数矩形用逐点枚举核对，并检查方程同乘负一、交换变量、特解平移、零系数、无解与64位端点构造。
这验证本地公式和建模，不新增模板、正式题覆盖或线上AC。
